\section{Original generic line sections and local orders}

Let $X$ be the SAME original integral scheme, $\eta$ its original generic
point, $K=\mathcal O_{X,\eta}$, and $\kappa_\eta$ its original residue
field.  An original generic section of the actual line $L$ means literally
\[
 \sigma\in\Gamma(\operatorname{Spec}\kappa_\eta,
             (\operatorname{fromSpecResidueField}\eta)^*L).
\]
It is not a supplied rational function or a new generic-line carrier.
The local-order construction below needs only integral locally Noetherian
$X$ and an actual invertible $L$.  Normality and projectivity will be used
by the subsequent global Euler argument, not by frame independence.

\begin{definition}[The actual generic coefficient in an original open frame]
\label{def:curve-generic-section-original-coefficient}
\leanok
\uses{def:actual-line-frame-via-factor,lem:actual-generic-residue-field-equivalence}
\lean{EnvelopeProperty.CurveGenericSectionOrders.genericFibreModule,
EnvelopeProperty.CurveGenericSectionOrders.genericCoefficient}
For a nonempty original open $U$, use the genuine factor
$q_U:\operatorname{Spec}\kappa_\eta\to U$ of the original residue-field
morphism.  For an actual frame $e:L|_U\simeq\mathcal O_U$, pull $e$ back by
this factor, using the original pullback-composition comparison.  Apply
this actual frame to $\sigma$, apply the native $\Gamma\operatorname{Spec}$
isomorphism, and then the inverse of the ORIGINAL residue map
$K\simeq\kappa_\eta$.  Denote the resulting element of $K$ by
$c_{U,e}(\sigma)$.  The module is named before its sections are formed,
with the same native scalar action and the same original pullback.
\end{definition}
\begin{proof}
Every nonempty original open contains the generic point.  The original
residue-field morphism has constant image $\eta$, so its genuine open
factor exists and its composite with the original inclusion is the
original morphism.  Pullback of an actual isomorphism and the actual
unit comparison give an isomorphism on the original residue fibre.
The original generic stalk is a field; its original residue quotient
is both injective and surjective.  Thus all maps in the stated
definition are the specified original maps and isomorphisms.
\end{proof}

\begin{lemma}[Faithful coefficients]
\label{lem:curve-generic-section-coefficient-restriction}
\leanok
\uses{def:curve-generic-section-original-coefficient,lem:actual-line-frame-via-factor-composition,scheme:actual-residue-section-pullback}
\lean{EnvelopeProperty.CurveGenericSectionOrders.genericCoefficient_ne_zero,
EnvelopeProperty.CurveGenericSectionOrders.genericCoefficient_restrict}
If $\sigma\ne0$, then $c_{U,e}(\sigma)\ne0$.  If $V\subseteq U$ is
nonempty and $e_V$ is the actual restriction of $e$ by the original
pullback-composition comparison, then
$c_{V,e_V}(\sigma)=c_{U,e}(\sigma)$.
\end{lemma}
\begin{proof}
An actual module-sheaf isomorphism gives inverse whole-section maps,
and both field isomorphisms are injective.  Hence the coefficient of a
nonzero original section is nonzero.  The two original factors
$q_V\to V\to U$ and $q_U\to U$ have equal composites with the open
immersion into $X$; they are therefore equal.  The already proved
whole-isomorphism composition law for actual frame pullbacks now
identifies the two frames, and hence every actual coefficient.
\end{proof}

\begin{lemma}[Original frame changes preserve the native local order]
\label{lem:curve-generic-section-order-frame-independent}
\leanok
\uses{lem:curve-generic-section-coefficient-restriction,lem:actual-unit-module-automorphism-scalar,def:fc-original-section-whole-frame-scalar}
\lean{EnvelopeProperty.CurveGenericSectionOrders.ord_genericCoefficient_eq}
Assume $X$ is locally Noetherian.  If $x\in U\cap V$, $e$ and $e'$ are
ANY two actual original open frames of $L$, and $\sigma\ne0$, then
\[
 \operatorname{ord}_x c_{U,e}(\sigma)
   =\operatorname{ord}_x c_{V,e'}(\sigma).
\]
\end{lemma}
\begin{proof}
Restrict both original frames to the actual intersection $W=U\cap V$,
which contains $x$.  Their actual composite is an automorphism of the
original structure module on $W$, hence multiplication by a genuine
original unit $v\in\Gamma(W,\mathcal O_W)^\times$.  Pulling this whole
scalar identity to the original generic point shows that the two
generic coefficients differ by the original generic germ of $v$.
Indeed, the native coefficient square for the residue-field morphism
is its original residue evaluation, followed by the inverse original
generic residue map, so this scalar is exactly that original germ.
The germ of $v$ is nonzero and has native order zero at $x$ because
$v$ is a unit on a neighborhood of $x$.  Native multiplicativity of
order for nonzero field elements, and the preceding restriction law,
give the claimed equality.  This holds at every original point,
including the generic point, without a DVR or order certificate.
\end{proof}

\begin{definition}[The intrinsic order of the original generic section]
\label{def:curve-generic-section-original-order}
\leanok
\uses{lem:curve-generic-section-order-frame-independent}
\lean{EnvelopeProperty.CurveGenericSectionOrders.order,
EnvelopeProperty.CurveGenericSectionOrders.order_eq_frameCoefficient,
EnvelopeProperty.CurveGenericSectionOrders.order_generic}
For $L$ invertible choose an actual frame neighborhood of each original
$x$, and define $\operatorname{ord}_x(\sigma)$ to be the native order of
the coefficient just defined.  For every nonzero original $\sigma$,
EVERY actual original frame $e$ on EVERY original $U\ni x$ computes it:
\[
 \operatorname{ord}_x(\sigma)=\operatorname{ord}_x c_{U,e}(\sigma).
\]
Its value at $\eta$ is zero.  The choice of frame is not a matching
assumption in this definition or theorem.
\end{definition}
\begin{proof}
The original invertibility atlas covers every original point, so
the chosen frame really exists.  The preceding frame independence
identifies its coefficient order with that of any given actual frame.
The original generic point has coheight zero, so its native order is
zero by the definition of native order away from coheight one.
\end{proof}

\begin{lemma}[Actual scalar and tensor section arithmetic]
\label{lem:curve-generic-section-scalar-tensor-order}
\leanok
\uses{def:curve-generic-section-original-order,def:actual-line-frame-via-factor,lem:actual-line-frame-via-factor-composition}
\lean{EnvelopeProperty.CurveGenericSectionOrders.scalarSection,
EnvelopeProperty.CurveGenericSectionOrders.genericCoefficient_scalarSection,
EnvelopeProperty.CurveGenericSectionOrders.exists_scalarSection,
EnvelopeProperty.CurveGenericSectionOrders.tensorSection,
EnvelopeProperty.CurveGenericSectionOrders.order_scalarSection,
EnvelopeProperty.CurveGenericSectionOrders.order_tensorSection}
The original isomorphism $K\simeq\kappa_\eta$ and native
$\Gamma\operatorname{Spec}$ inverse make the actual scalar multiple
$f\sigma$ in the SAME original generic fibre.  Its coefficient is
$f c_{U,e}(\sigma)$.  Every two nonzero original sections of the same
generic line differ by ONE nonzero scalar $f\in K$.
For $\sigma\ne0$ and $\tau\ne0$, the actual generic pure tensor,
formed with the original pullback tensorator, has local order equal to
the sum of the two factor orders.  For $f\ne0$ and $\sigma\ne0$, scalar
multiplication adds $\operatorname{ord}_x f$ to section order.
\end{lemma}
\begin{proof}
The actual whole-section map of a frame is linear over the original
residue-spectrum section ring.  The two actual ring isomorphisms
therefore give the scalar formula.  In any ONE original generic
frame, set $f$ equal to the quotient of the two nonzero coefficients.
The scalar formula and injectivity of that actual frame prove equality
of the original sections.  Thus this is one global scalar, not a family
of pointwise coefficient matches.

On a common original open, tensor the two original frames using the
native pullback tensorator and the original structure-line tensor
comparison.  The monoidal pullback-composition com\-parison identifies
the resulting actual ge\-neric frame with the tensor of the two actual
generic frames.  The original pure-tensor section is mapped to the
product of their original coefficients.  Both coefficients are
nonzero, so this also proves the original tensor section is nonzero.
Native multiplicativity of order and the ANY-frame theorem give the
stated scalar and tensor order formulas.
\end{proof}

\begin{lemma}[Original whole sections retain their actual zero orders]
\label{lem:curve-generic-section-whole-zero-order}
\leanok
\uses{def:curve-generic-section-original-order,def:fc-original-line-section-frame-coefficient,thm:fc-original-line-section-frame-order,lem:fc-original-local-section-scalar,lem:scheme-original-global-section-hom,lem:fc-original-line-section-fibre,thm:line-section-zero-ideal-cartier,def:fc-curve-cartier-original-equations}
\lean{EnvelopeProperty.CurveGenericSectionOrders.wholeSection,
EnvelopeProperty.CurveGenericSectionOrders.wholeSection_ne_zero,
EnvelopeProperty.CurveGenericSectionOrders.order_wholeSection}
For every nonzero ORIGINAL whole section $s:\mathcal O_X\to L$, its
actual original generic pullback $s_\eta$ is nonzero, and for every
original $x$ its intrinsic order is the old native local order of the
SAME original zero ideal of $s$.
\end{lemma}
\begin{proof}
Choose a genuine original affine frame at $x$.  The already proved
whole local section-Hom identity says that the original generic
coefficient is exactly the generic germ of the actual local scalar.
The native over-site frame and native open-subscheme frame coefficient
comparison identify that scalar with the SAME original zero-ideal
coefficient.  The original nonzero-section theorem proves this
coefficient is nonzero on every nonempty such open.  Hence the
original generic section cannot be zero.  The already proved
zero-ideal local-order formula and ANY-frame theorem give the claimed
equality at every original point.
\end{proof}

\begin{theorem}[One global rational correction of two actual whole zero divisors]
\label{thm:curve-generic-section-generated-order-difference}
\leanok
\uses{lem:curve-generic-section-scalar-tensor-order,lem:curve-generic-section-whole-zero-order,thm:fc-original-projective-generated-line-quotients,thm:fc-original-line-nonzero-generator,thm:line-section-zero-ideal-cartier}
\lean{EnvelopeProperty.CurveGenericSectionOrders.exists_order_difference}
On the SAME original integral projective $X$, for every actual
invertible $L$ and every nonzero original generic $\sigma$, there are
an internally produced invertible denominator $T$, nonzero original
WHOLE sections $s$ of $L\otimes T$ and $t$ of $T$, and ONE nonzero
$f\in K$, such that at EVERY original point
\[
 \operatorname{ord}_x\sigma
  =\operatorname{ord}_x f+
    \operatorname{ord}_x(s)-\operatorname{ord}_x(t).
\]
The two whole-section orders are literally the old local orders of
their actual original zero ideals.  No coefficient, order, divisor,
finite-support, or Euler comparison is supplied.
\end{theorem}
\begin{proof}
Actual projective generation supplies one genuine denominator $T$
such that both $T$ and $L\otimes T$ have genuine global generating
families.  The original nonzero-generator theorem internally chooses
nonzero WHOLE $t$ and $s$.  Their actual generic pullbacks are nonzero.
The true generic tensor $\sigma\otimes t_\eta$ is also nonzero.
The preceding one-dimensional original-section argument gives ONE
nonzero global scalar $f$ with
$\sigma\otimes t_\eta=f s_\eta$.  Compute the orders of both sides in
ANY original local frames, apply the true tensor and scalar formulas,
and substitute the two actual zero-ideal orders.  Rearrangement gives
the stated formula simultaneously for all points, including poles.
\end{proof}

\begin{theorem}[Arbitrary original generic sections compute the SAME full Euler degree]
\label{thm:curve-generic-section-full-euler-degree}
\leanok
\uses{thm:curve-generic-section-generated-order-difference,def:fc-original-weighted-rational-order,thm:fc-original-curve-principal-weighted-balance,thm:fc-original-curve-line-degree-zero-section,thm:fc-original-curve-line-degree-tensor,thm:fc-curve-cartier-ideal-finite,thm:fc-curve-finite-closed-subschemes,thm:fc-curve-cartier-native-local-orders,def:fc-curve-cartier-original-equations,thm:line-section-zero-ideal-cartier}
\lean{EnvelopeProperty.CurveGenericSectionOrders.hasFiniteSupport_order,
EnvelopeProperty.CurveGenericSectionOrders.degree_eq_finsum_order}
Let $p:X\to\operatorname{Spec}k$ be the SAME original projective
integral curve, $\dim X\le1$, with every original local ring normal.
Then EVERY nonzero ORIGINAL generic section $\sigma$ of EVERY
invertible original $L$ has finite-support intrinsic local order and
\[
 \deg_p L=\sum_{x\in X}[\kappa(x):k]\operatorname{ord}_x\sigma.
\]
The degree is the original normalized FULL coherent Euler
characteristic.

Residue scalar actions are the actual ones induced
by $p$.  The sum is over ALL original points and keeps negative poles.
\end{theorem}
\begin{proof}
Projectivity derives original Noetherianity and compactness.  The one
global correction formula expresses the section orders as the sum
and difference of three honest finite-support functions: the order
of ONE nonzero global rational function (finite by normal-curve
Cartier denominator geometry), and the two zero orders (supported
on their actual proper Cartier supports, which are finite on the
original curve).  This proves finite support internally.

Multiply that same pointwise identity by each true residue degree.
The weighted sum of the ONE global rational-function orders is zero
by the genuine original principal-divisor balance.  The old TWO
nonzero whole-section Euler theorems give the weighted zero orders
as $\deg_p(L\otimes T)$ and $\deg_p T$.  Actual full Euler tensor
additivity gives
$\deg_p(L\otimes T)-\deg_p T=\deg_pL$.  Finite-support sum
additivity therefore proves exactly the stated original degree
formula.  No original vertical ideal is restricted to the curve,
and no generic section is treated as a whole global section.
\end{proof}
